% ═══════════════════════════════════════════════════════════════
% SLIDES: Describir fotos
% Klasse 8 · Spanisch · Beamer
% ═══════════════════════════════════════════════════════════════

\documentclass[aspectratio=169, 14pt]{beamer}

\usepackage[utf8]{inputenc}
\usepackage[T1]{fontenc}
\usepackage[ngerman]{babel}
\usepackage{palatino}
\usepackage{tikz}
\usepackage{booktabs}
\usepackage{array}

% ═══════════════════════════════════════════════════════════════
% THEME
% ═══════════════════════════════════════════════════════════════

\usetheme{default}
\usecolortheme{default}

\definecolor{spanisch}{HTML}{C62828}
\definecolor{akzent}{HTML}{1565C0}
\definecolor{success}{HTML}{2E7D32}
\definecolor{hellgrau}{HTML}{F5F5F5}

\setbeamercolor{structure}{fg=spanisch}
\setbeamercolor{frametitle}{fg=spanisch}
\setbeamercolor{title}{fg=white, bg=spanisch}
\setbeamercolor{block title}{fg=white, bg=spanisch}
\setbeamercolor{block body}{bg=hellgrau}

\setbeamertemplate{navigation symbols}{}
\setbeamertemplate{footline}[frame number]

% ═══════════════════════════════════════════════════════════════
% BEFEHLE
% ═══════════════════════════════════════════════════════════════

\newcommand{\es}[1]{\textbf{#1}}
\newcommand{\de}[1]{\textit{\textcolor{gray}{#1}}}
\newcommand{\abmark}[1]{\textcolor{spanisch}{$\rightarrow$ AB #1}}

% ═══════════════════════════════════════════════════════════════
% TITELFOLIE
% ═══════════════════════════════════════════════════════════════

\title{Describir fotos}
\subtitle{Bildbeschreibung auf Spanisch}
\author{Spanisch · Klasse 8}
\date{}

\begin{document}

% --- F1: Titel ---
\begin{frame}
    \titlepage
\end{frame}

% --- F2: Einstieg ---
\begin{frame}{¿Qué ves en la foto?}

\begin{center}
    \Large
    Stell dir vor, du zeigst einem spanischen Freund\\
    ein Foto aus deinem Urlaub.\\[1em]

    \textbf{Wie beschreibst du, was du siehst?}
\end{center}

\vspace{1em}

\begin{block}{Heute lernst du:}
\begin{itemize}
    \item Aktionen beschreiben (presente continuo)
    \item Bilder strukturiert beschreiben
    \item Positionen und Personen beschreiben
\end{itemize}
\end{block}

\end{frame}

% --- F3: Englisch-Brücke ---
\begin{frame}{Englisch-Brücke: Present Progressive}

\begin{center}
\Large

\textcolor{akzent}{I \textbf{am reading}} = \textcolor{spanisch}{\es{Estoy leyendo}}

\vspace{1.5em}

\textcolor{akzent}{She \textbf{is walking}} = \textcolor{spanisch}{\es{Ella está caminando}}

\vspace{1.5em}

\textcolor{akzent}{They \textbf{are playing}} = \textcolor{spanisch}{\es{Ellos están jugando}}

\end{center}

\vspace{1em}

\begin{block}{}
\centering
Gleiche Struktur: \textbf{to be + -ing} = \textbf{estar + gerundio}
\end{block}

\end{frame}

% --- F4: Presente Continuo Regel ---
\begin{frame}{El presente continuo}

\begin{block}{Bildung}
\centering
\Large
\es{estar} (konjugiert) + \es{gerundio}
\end{block}

\vspace{1em}

\begin{columns}
\begin{column}{0.45\textwidth}
\textbf{estar:}
\begin{tabular}{ll}
yo & estoy \\
tú & estás \\
él/ella & está \\
nosotros & estamos \\
vosotros & estáis \\
ellos & están \\
\end{tabular}
\end{column}

\begin{column}{0.5\textwidth}
\textbf{Gerundio-Endungen:}

\vspace{0.5em}

\es{-ar} $\rightarrow$ \es{-ando}\\
\de{hablar → hablando}

\vspace{0.5em}

\es{-er/-ir} $\rightarrow$ \es{-iendo}\\
\de{comer → comiendo}
\end{column}
\end{columns}

\vspace{1em}
\abmark{G1 ausfüllen}

\end{frame}

% --- F5: Gerundio Beispiele ---
\begin{frame}{Verbos en gerundio}

\begin{columns}
\begin{column}{0.5\textwidth}
\begin{tabular}{ll}
\toprule
\textbf{Infinitivo} & \textbf{Gerundio} \\
\midrule
caminar & caminando \\
pasear & paseando \\
mirar & mirando \\
hablar & hablando \\
comer & comiendo \\
beber & bebiendo \\
\bottomrule
\end{tabular}
\end{column}

\begin{column}{0.5\textwidth}
\begin{tabular}{ll}
\toprule
\textbf{Infinitivo} & \textbf{Gerundio} \\
\midrule
leer & \textcolor{spanisch}{leyendo} \\
escribir & escribiendo \\
jugar & jugando \\
correr & corriendo \\
reír & \textcolor{spanisch}{riendo} \\
dormir & durmiendo \\
\bottomrule
\end{tabular}
\end{column}
\end{columns}

\vspace{1em}

\begin{alertblock}{Achtung}
\es{leer → leyendo} (nicht "leiendo") \quad \es{oír → oyendo}
\end{alertblock}

\abmark{A1: Alle 12 Verben}

\end{frame}

% --- F6: Bildaufbau ---
\begin{frame}{Las partes de la imagen – Bildaufbau}

\begin{center}
\begin{tikzpicture}[scale=0.8]
    % Bild-Rahmen
    \draw[thick] (0,0) rectangle (12,7);

    % Bereiche
    \fill[yellow!20] (0,0) rectangle (12,2.5);
    \fill[orange!10] (0,2.5) rectangle (12,4.5);
    \fill[blue!10] (0,4.5) rectangle (12,7);

    % Labels
    \node at (6,1.25) {\Large \es{en primer plano} \de{(Vordergrund)}};
    \node at (6,3.5) {\Large \es{en segundo plano} \de{(Mittelgrund)}};
    \node at (6,5.75) {\Large \es{al fondo} \de{(Hintergrund)}};

    % Seiten
    \node[rotate=90] at (-0.8,3.5) {\es{a la izquierda}};
    \node[rotate=-90] at (12.8,3.5) {\es{a la derecha}};

    % Mitte
    \draw[dashed, spanisch] (6,0) -- (6,7);
    \node at (6,-0.5) {\es{en el centro}};
\end{tikzpicture}
\end{center}

\end{frame}

% --- F7: Bildaufbau Vokabeln ---
\begin{frame}{Las partes de la imagen – Vokabeln}

\begin{columns}
\begin{column}{0.5\textwidth}
\begin{tabular}{ll}
\es{en primer plano} & im Vordergrund \\[0.5em]
\es{en segundo plano} & im Mittelgrund \\[0.5em]
\es{al fondo} & im Hintergrund \\[0.5em]
\end{tabular}
\end{column}

\begin{column}{0.5\textwidth}
\begin{tabular}{ll}
\es{a la izquierda} & links \\[0.5em]
\es{a la derecha} & rechts \\[0.5em]
\es{en el centro} & in der Mitte \\[0.5em]
\end{tabular}
\end{column}
\end{columns}

\vspace{1.5em}

\begin{exampleblock}{Beispiele}
\es{En primer plano hay un banco.}\\
\es{Al fondo se ven edificios.}\\
\es{A la derecha, una mujer está caminando.}
\end{exampleblock}

\abmark{A2: Zuordnung}

\end{frame}

% --- F8: Positionen ---
\begin{frame}{Las posiciones – Positionen}

\begin{center}
\Large

\begin{tabular}{ll}
\es{estar de pie} & stehen \\[0.8em]
\es{estar sentado/a} & sitzen \\[0.8em]
\es{estar tumbado/a} & liegen \\[0.8em]
\es{estar apoyado/a} & sich anlehnen \\
\end{tabular}

\end{center}

\vspace{1em}

\begin{alertblock}{Achtung: Endung passt zum Geschlecht!}
\es{El hombre está sentado.} / \es{La mujer está sentada.}
\end{alertblock}

\abmark{A3: Lückentext}

\end{frame}

% --- F9: Personenbeschreibung Repaso ---
\begin{frame}{Describir personas – Repaso}

\begin{columns}
\begin{column}{0.5\textwidth}
\textbf{Größe:}
\begin{itemize}
    \item \es{alto/a} – groß
    \item \es{bajo/a} – klein
\end{itemize}

\vspace{0.5em}

\textbf{Haare:}
\begin{itemize}
    \item \es{el pelo largo/corto}
    \item \es{rubio/moreno/negro}
    \item \es{liso/rizado}
\end{itemize}
\end{column}

\begin{column}{0.5\textwidth}
\textbf{Augen:}
\begin{itemize}
    \item \es{los ojos azules}
    \item \es{verdes/marrones}
\end{itemize}

\vspace{0.5em}

\textbf{Sonstiges:}
\begin{itemize}
    \item \es{llevar gafas}
    \item \es{joven/mayor}
\end{itemize}
\end{column}
\end{columns}

\end{frame}

% --- F10: Orte Park/Stadt ---
\begin{frame}{Los lugares – Orte}

\begin{columns}
\begin{column}{0.5\textwidth}
\textbf{Im Park:}
\begin{tabular}{ll}
\es{el parque} & der Park \\
\es{el banco} & die Bank \\
\es{el árbol} & der Baum \\
\es{el césped} & der Rasen \\
\es{la fuente} & der Brunnen \\
\es{el camino} & der Weg \\
\end{tabular}
\end{column}

\begin{column}{0.5\textwidth}
\textbf{In der Stadt:}
\begin{tabular}{ll}
\es{la calle} & die Straße \\
\es{la plaza} & der Platz \\
\es{el edificio} & das Gebäude \\
\es{la tienda} & das Geschäft \\
\end{tabular}

\vspace{0.5em}

\textbf{Im Café:}
\begin{tabular}{ll}
\es{la cafetería} & das Café \\
\es{la terraza} & die Terrasse \\
\end{tabular}
\end{column}
\end{columns}

\end{frame}

% --- F11: Redemittel ---
\begin{frame}{Redemittel – Satzanfänge}

\begin{block}{Einstieg}
\es{En la foto veo...} \de{(Auf dem Foto sehe ich...)}\\
\es{La imagen muestra...} \de{(Das Bild zeigt...)}
\end{block}

\begin{block}{Bildaufbau}
\es{En primer plano hay...} \de{(Im Vordergrund gibt es...)}\\
\es{Al fondo se ve...} \de{(Im Hintergrund sieht man...)}
\end{block}

\begin{block}{Aktionen + Vermutungen}
\es{Una persona está + gerundio...}\\
\es{Parece que...} \de{(Es scheint, dass...)}
\end{block}

\end{frame}

% --- F12: Jetzt du! ---
\begin{frame}{¡Ahora tú! – Jetzt bist du dran!}

\begin{center}
\Large

Beschreibe das Foto auf deinem Arbeitsblatt!

\vspace{1.5em}

\normalsize
\textbf{Nutze:}
\begin{itemize}
    \item Bildaufbau (en primer plano, al fondo...)
    \item Presente continuo (estar + gerundio)
    \item Positionen (de pie, sentado, tumbado)
    \item Personenbeschreibung (pelo, ojos, alto)
    \item Redemittel (En la foto veo...)
\end{itemize}

\vspace{1em}

\abmark{A4: Analyse + A5: Eigener Text}

\end{center}

\end{frame}

% --- F13: Checkliste ---
\begin{frame}{Checkliste: Gute Bildbeschreibung}

\begin{center}
\Large

\begin{itemize}
    \item[$\square$] Bildaufbau verwendet
    \item[$\square$] Presente continuo korrekt
    \item[$\square$] Positionen beschrieben
    \item[$\square$] Personen beschrieben
    \item[$\square$] Orte benannt
    \item[$\square$] Redemittel genutzt
\end{itemize}

\vspace{1em}

\normalsize
\abmark{Checkliste K1}

\end{center}

\end{frame}

\end{document}
